\documentclass[11pt]{article}

% --- Packages ---------------------------------------------------------------
\usepackage[margin=0.75in,top=0.7in,bottom=0.7in]{geometry}
\usepackage{amsmath}
\usepackage{microtype}
\usepackage{booktabs}
\usepackage{array}
\usepackage{hyperref}
\usepackage{xcolor}
\usepackage{listings}
\usepackage{enumitem}
\usepackage{titlesec}
\usepackage{caption}

\hypersetup{colorlinks=true,linkcolor=black,urlcolor=blue!50!black,citecolor=blue!50!black}

\titleformat{\section}{\large\bfseries}{\thesection.}{0.5em}{}
\titleformat{\subsection}{\normalsize\bfseries}{\thesubsection}{0.5em}{}
\titlespacing*{\section}{0pt}{0.7ex}{0.3ex}
\titlespacing*{\subsection}{0pt}{0.4ex}{0.2ex}
\setlist{nosep,leftmargin=1.4em,topsep=1pt}
\setlength{\parskip}{2pt plus 0.5pt}

\lstset{basicstyle=\ttfamily\footnotesize,breaklines=true,columns=fullflexible,showstringspaces=false}

% --- Title ------------------------------------------------------------------
\title{\vspace{-2em}\textbf{Project 2: Schema Linking with Supervised Fine-Tuning}\\
\large CSE/DSC 234 -- Data-Centric AI and AI Engineering, Spring 2026}
\author{Jay Chaudhary \quad\textbar\quad Abhay Jain}
\date{June 4, 2026}

\begin{document}
\maketitle
\vspace{-2.8em}

% ===========================================================================
\section{Training Data Methodology}
% ===========================================================================

\paragraph{Base split.}
We use the released 301 training examples and 101 validation examples from the SNAILS-derived
dataset without any external augmentation. The 17 database schemas span national-parks
biodiversity, NTSB crash statistics, NY education data, and SAP business-resource planning.
We deliberately chose \emph{not} to augment with Spider or synthesised NL-SQL pairs because
schema serialisation and prompt design proved to be higher-leverage axes: fixing a broken schema
format moved the score by up to 0.4 points, whereas preliminary augmentation runs introduced
noise from domain mismatch.

\paragraph{Schema serialisation.}
Each schema is rendered as one line per table in the format:
\begin{lstlisting}
TABLE(col*, col>RefTable, col, ...)
\end{lstlisting}
where \texttt{*} marks a primary key and \texttt{>RefTable} marks a foreign key.
No spaces after commas. This encoding fits all relational structure (joins, PK identity) in
minimal tokens, keeping 94\% of training examples below 2048 tokens without truncation.
Serialisation is implemented in \texttt{format\_data.py::serialize\_compact}, reading
\texttt{primary\_keys} and \texttt{foreign\_keys} directly from Spider-format schema JSON.

\paragraph{Output format.}
The gold target is the schema-links dict serialised as compact JSON
(\texttt{separators=(",",":")}).  Tables referenced without specific columns (e.g.\
\texttt{COUNT(*)}) are emitted as \texttt{\{"T":[]\}}, matching the wildcard-boundary rule
in the eval script exactly.

\paragraph{Chat record construction.}
Each example becomes a three-turn supervised chat record:
\textit{system} (35-token compact instruction),
\textit{user} (question first, then \texttt{DB:} and \texttt{Schema:} block),
\textit{assistant} (compact JSON).
Question-first ordering was motivated by the hypothesis that the decoder should attend to
question tokens while scanning each schema line; ablation confirmed a +0.09 score gain.
Records are saved as HuggingFace \texttt{Dataset} Arrow files via \texttt{format\_data.py}.

% ===========================================================================
\section{Final Pipeline Architecture}
% ===========================================================================

Table~\ref{tab:final} summarises the best configuration (C8, score \textbf{0.6117}).

\begin{table}[h]
\centering\small
\caption{Final pipeline configuration (C8).}
\label{tab:final}
\begin{tabular}{l l}
\toprule
\textbf{Component} & \textbf{Setting} \\
\midrule
Base model          & \texttt{Qwen/Qwen2.5-1.5B-Instruct} ($\le$2B constraint) \\
LoRA targets        & \texttt{q\_proj, k\_proj, v\_proj, o\_proj} \\
LoRA rank / alpha   & $r=64$, $\alpha=128$, dropout $=0.1$ \\
Learning rate       & $5\times10^{-5}$, cosine schedule, 5\% linear warmup \\
Batch / accumulation & 1 per device $\times$ 4 gradient-accumulation steps (eff.\ batch 4) \\
Epochs / seq.\ length & 5 epochs, max 2048 tokens \\
Precision           & \texttt{bf16} mixed precision \\
Schema format       & \texttt{TABLE(col*,col>Ref)} compact; PK=\texttt{*}, FK=\texttt{>Ref} \\
Prompt layout       & Question first: \texttt{Q:\{q\}} then \texttt{DB:\{id\}}, \texttt{Schema:}\{schema\} \\
System prompt       & 35-token compact instruction (no verbose preamble) \\
Decoding            & Greedy (\texttt{do\_sample=False}), repetition penalty 1.1 \\
Post-processing     & \texttt{filter\_to\_schema}: case-insensitive hallucination removal \\
\bottomrule
\end{tabular}
\end{table}

\paragraph{Fine-tuning (\texttt{train.py}).}
LoRA adapters are applied via RapidFire AI's \texttt{RFLoraConfig} and trained with
\texttt{RFSFTConfig} wrapping HuggingFace TRL's \texttt{SFTTrainer}.  Experiments are
managed through \texttt{Experiment} + \texttt{RFGridSearch}; each config runs to completion
(\texttt{num\_chunks=1}) before checkpoints are compared by validation loss.  The best
checkpoint is copied to \texttt{./adapter/}.

\paragraph{Inference (\texttt{main.py}).}
The base model is loaded from HuggingFace Hub; the LoRA adapter is attached via
\texttt{PeftModel.from\_pretrained(base, "./adapter")}.  Raw model output is parsed by a
tolerant extractor: direct \texttt{json.loads}, then regex on the first \texttt{\{\ldots\}}
block, then empty-dict fallback.  \texttt{filter\_to\_schema} removes any hallucinated
table or column not present in the actual schema, using case-insensitive lookup maps to
restore exact casing.  Both adapter location and base model ID are auto-detected from
\texttt{adapter/adapter\_config.json} so the script is self-contained.

% ===========================================================================
\section{Experimentation Methodology}
\label{sec:exp}
% ===========================================================================

We ran all configs through RapidFire AI's \texttt{Experiment} / \texttt{RFGridSearch} API,
using \texttt{RFLoraConfig}, \texttt{RFModelConfig}, and \texttt{RFSFTConfig} to define each
knob bundle.  Configs were compared by validation loss logged to \texttt{logs/} and by
end-to-end leaderboard score on the 101-question validation set.  Underperforming configs were
stopped via the IC Ops interface; promising ones were cloned and modified.

Table~\ref{tab:configs} lists all 9 distinct configurations explored across five knob axes:
schema format, learning rate, LoRA rank, prompt layout, and system-prompt wording.

\begin{table}[h]
\centering\footnotesize
\caption{All configurations explored. ``Meas.''\ = full train + eval on validation set.}
\label{tab:configs}
\setlength{\tabcolsep}{3.5pt}
\begin{tabular}{>{\raggedright\arraybackslash}p{2.5cm} c c c >{\raggedright\arraybackslash}p{3.8cm} c}
\toprule
\textbf{Config} & $r$ & \textbf{LR} & \textbf{Ep.} & \textbf{Schema / Prompt} & \textbf{Status} \\
\midrule
C1 Baseline       & 16  & $2\times10^{-4}$ & 3 & Simple \texttt{TABLE(c1,c2)}, schema-first  & Meas. \\
C2 high LR        & 64  & $1\times10^{-4}$ & 5 & Simple, schema-first                         & Meas. \\
C3 low LR         & 64  & $8\times10^{-6}$ & 5 & Simple, schema-first                         & Meas. \\
C4 JSON hier.     & 64  & $5\times10^{-5}$ & 5 & Nested JSON with types, PK, FK               & Meas. \\
C5 rich inline    & 64  & $5\times10^{-5}$ & 5 & \texttt{TABLE(col [PK,TEXT],\ldots)}          & Meas. \\
C6 r64 simple     & 64  & $5\times10^{-5}$ & 5 & Simple, schema-first                         & Meas. \\
C7 compact+prune  & 64  & $5\times10^{-5}$ & 5 & Compact PK/FK, Q-first, keyword pruning      & Meas. \\
\textbf{C8 compact (best)} & 64 & $5\times10^{-5}$ & 5 & Compact PK/FK, Q-first, \textbf{no prune}  & \textbf{Meas.} \\
C9 r128+col prompt & 128 & $5\times10^{-5}$ & 5 & Compact PK/FK, Q-first, col-aware sys.~prompt & Meas. \\
\bottomrule
\end{tabular}
\end{table}

\paragraph{Learning rate search (C1--C3, C6).}
With 301 training examples and effective batch 4 (75 steps/epoch), LR is the most sensitive
knob. At $\ge10^{-4}$ (C1, C2) token accuracy oscillated between 0.42 and 0.56: the model
peaked near step~50 (epoch~1) then degraded — a classic small-dataset overshoot.
At $8\times10^{-6}$ (C3) convergence was too slow (score 0.47).
$5\times10^{-5}$ with cosine schedule and 5\% warmup gave stable, monotonically improving
validation loss across all 5 epochs.

\paragraph{Schema format search (C4--C8).}
Nested JSON (C4) frequently exceeded 2048 tokens on databases with 20+ tables, causing
truncation (score 0.20). Rich inline annotation (C5) added too many tokens and confused
the model (0.45). The simple \texttt{TABLE(col1,col2)} format (C6) reached 0.54.
The compact PK/FK variant (C7, C8) encodes relational structure in near-zero extra tokens:
\texttt{col*} for PKs and \texttt{col>Ref} for FKs tells the model which columns are join
keys without verbose type strings.

\paragraph{Pruning ablation (C7 vs.\ C8).}
C7 pruned the schema at inference time to keyword-matching tables plus FK neighbours.
Despite reducing median context length, table score dropped from 0.683 to 0.552 because
pruning has no oracle knowledge at inference -- tables whose names share no lexical tokens
with the question (e.g., a lookup table referenced only through a FK chain) were silently
hidden. Removing pruning (C8) recovered the table score and set the best overall result.

\paragraph{Rank and column-aware prompt (C9).}
Doubling the rank to $r=128$ and adding an explicit column-type enumeration to the system
prompt (\textit{``List ALL columns: SELECT outputs, WHERE/HAVING filters, JOIN keys, GROUP BY
and ORDER BY columns''}) targeted the column recall gap. Score was 0.6074 -- marginally
below C8 -- suggesting the column gap is limited by data scale rather than adapter capacity
or instruction wording at this stage.

% ===========================================================================
\section{Results and Analysis}
% ===========================================================================

\begin{table}[h]
\centering\footnotesize
\caption{Results on the 101-question validation set.}
\label{tab:results}
\setlength{\tabcolsep}{5pt}
\begin{tabular}{l c c c c c c}
\toprule
\textbf{Config} & $r$ & \textbf{LR} & \textbf{Ep.} & \textbf{Table Score} & \textbf{Column Score} & \textbf{Score} \\
\midrule
C1 Baseline    & 16  & $2\times10^{-4}$ & 3 & --    & --    & 0.520 \\
C2 high LR     & 64  & $1\times10^{-4}$ & 5 & --    & --    & 0.520 \\
C3 low LR      & 64  & $8\times10^{-6}$ & 5 & --    & --    & 0.470 \\
C4 JSON hier.  & 64  & $5\times10^{-5}$ & 5 & --    & --    & 0.200 \\
C5 rich inline & 64  & $5\times10^{-5}$ & 5 & --    & --    & 0.450 \\
C6 r64 simple  & 64  & $5\times10^{-5}$ & 5 & --    & --    & 0.540 \\
C7 +prune      & 64  & $5\times10^{-5}$ & 5 & 0.552 & 0.484 & 0.518 \\
\textbf{C8 (best)} & \textbf{64} & $5\times10^{-5}$ & \textbf{5} & \textbf{0.683} & \textbf{0.541} & \textbf{0.612} \\
C9 r128+col    & 128 & $5\times10^{-5}$ & 5 & 0.685 & 0.529 & 0.607 \\
\bottomrule
\end{tabular}
\end{table}

\paragraph{What worked.}
The compact PK/FK schema format combined with question-first prompt layout produced the
largest single gain ($+0.07$ over simple format). Embedding relational structure inline
improved both table and column scores: the model can infer join paths from \texttt{>Ref}
markers without a separate ``Relationships'' section. Stable learning rate ($5\times10^{-5}$)
was essential on this data scale -- higher rates caused oscillation and lower rates failed to
converge in 5 epochs.

\paragraph{What didn't work.}
Richer schema formats (C4, C5) hurt rather than helped: JSON hierarchical truncated long
schemas, and rich inline annotations added noise tokens without proportional signal gain.
Inference-time schema pruning (C7) degraded table recall by hiding tables that had no
lexical overlap with the question. Larger adapter rank (C9) and an explicit column-type
instruction in the system prompt did not improve column recall, suggesting the gap is
data-scale limited.

\paragraph{Table vs.\ column gap.}
The table score (0.683) exceeds the column score (0.541) by 0.14 throughout all runs.
Column prediction is harder: the model must identify implicit references -- join keys,
WHERE predicates, GROUP BY columns -- that are not lexically present in the question.
Zero schema-invalid identifiers were predicted in C8 and C9, confirming that
\texttt{filter\_to\_schema} eliminates all hallucinations.

\paragraph{Surprises.}
Removing schema pruning (C8 vs.\ C7) \emph{increased} score despite longer contexts.
The 2048-token limit was not the bottleneck; hiding relevant tables was.
The column-aware system prompt (C9) was expected to lift column recall but had no effect,
implying that the model already ``knows'' to look for filter and join columns -- the problem
is that the small training set does not cover enough distinct join patterns to generalise.

\paragraph{With more time/budget.}
(1) Augment with LLM-generated NL-SQL pairs labelled by \texttt{sql\_to\_schema\_links.py},
targeting under-represented SAP sub-modules ($\sim$6 training examples each).
(2) Curriculum training: start with single-table questions, progress to multi-table joins.
(3) Sorted output: always emit tables and columns in schema-defined order to reduce
decoding variance and improve $F_1$.
(4) Custom diagnostic metric: \textit{join-coverage score} -- fraction of FK-join columns
correctly predicted -- to target the implicit-reference gap directly.

% ===========================================================================
\section{AI Tools Statement}
% ===========================================================================

Claude Code (Anthropic) was used throughout this project as a coding assistant: generating
boilerplate for \texttt{format\_data.py}, \texttt{train.py}, and \texttt{main.py};
debugging schema-parsing index errors; iterating on schema serialisation formats; and
drafting this report. All experimental design decisions, hyperparameter choices, result
interpretation, and final report content are the responsibility of the author.
Training scripts, prompt designs, and pipeline decisions were reviewed, tested, and
validated end-to-end by the author on DSMLP before inclusion.

% ---------------------------------------------------------------------------
\bibliographystyle{plain}
\begin{thebibliography}{9}
\bibitem{lora}
E.~J.\ Hu et al.\ LoRA: Low-Rank Adaptation of Large Language Models. \textit{ICLR 2022}.
\bibitem{spider}
T.~Yu et al.\ Spider: A Large-Scale Human-Labeled Dataset for Complex and Cross-Domain
Semantic Parsing and Text-to-SQL Task. \textit{EMNLP 2018}.
\bibitem{snails}
K.~Luoma and A.~Kumar. SNAILS: Schema Naming Assessments for Improved LLM-Based SQL
Inference. \textit{ACM SIGMOD 2025}.
\end{thebibliography}

\end{document}
